% Section 8 (part 7): the key lemma, the residue, Theorem P(b).

\begin{lem}[The key lemma]\genusstatus{AI-proposed, unchecked}
\label{lem:con-key}
Expand $\Sigma:=\sum_b\pi_{\nu_b}\,s^{-\nu_b}\,\IC_{\nu_b}(v(u),z)$ into monomials $u^Ez^\beta$, writing $\zh=s^{-4}z$ in $\pi_m$,
\[
\IC_m(v,z)=\sum_{i'\geq0}\frac{(-1/(4z))^{i'}}{i'!}\,\frac{(-v/z)^e}{e!},\qquad e:=m+4i'+1\quad(\text{terms with }e\geq0),
\]
and $v^e=u^eV^e$ (\cref{def:gau-van}; $v(u)=u\,\tau^{-5/4}$ by \cref{def:set-tau}, and $V=\tau^{-5/4}$ by \cref{lem:con-cfu}). A monomial $u^i\zh^j$ of $\pi_m$, the $i'$-th term of $\IC_m$ and a monomial $u^l$ of $V^e$ contribute to $u^Ez^\beta$ with $\beta=j-5i'-m-1$ and $E=i+l-(4\beta-1)/5$. Hence only $\beta\equiv4\bmod5$ occurs, and the coefficient of $u^Ez^\beta$ receives contributions only from the $u^i$-coefficients of the $\pi_m$ with $i\leq E+(4\beta-1)/5$, through finitely many terms; $\Sigma$ is a well-defined formal sum of monomials $u^Ez^\beta$.
\end{lem}

\begin{proof}
The $z$-exponent is $j-i'-e=j-5i'-m-1$, and $j\equiv m\bmod5$ by \cref{lem:con-red}, so $\beta\equiv-1$. The $u$-exponent is $i-\frac{4j}5-\frac m5+e+l$; substituting $j=\beta+5i'+m+1$ and $e=m+4i'+1$ gives $E=i+l-(4\beta-1)/5$. So $i=E+(4\beta-1)/5-l\leq E+(4\beta-1)/5$; for fixed $(E,\beta)$ and $i$ the exponent $l$ is fixed, $\pi_m$ is a polynomial in $\zh$ at the order $u^i$, so finitely many $j$ occur, and then $i'$ is fixed.
\end{proof}

\begin{prop}[The residue]\genusstatus{AI-proposed, unchecked}
\label{prop:con-res}
Work in a conifold frame, and let $P$ be a graded regular gauge with coefficients in $\bbQ\fps u$ with $\Phi_P=\PhiSP$ (\cref{lem:con-id}). Then
\[
\cV=\varpi_0\,\bigl(P\,\YC(v(u),z)\bigr)_0\qquad\text{in }\bbQ\lser{z}\fps{u}.
\]
\end{prop}

\begin{proof}
By \cref{lem:con-relab}, $\cV=s\hat\cV$ coefficientwise. By \cref{lem:fsp-res} the residue on $\cR_{\mathrm{res}}\fps{u}$ is invariant under the uniformising map, termwise in $u$, and vanishes on exact forms; with \cref{lem:con-red},
\[
\hat\cV=\Res_Y\bigl[\re^{-W_V/\zh}\fa_Y\,\dd Y\bigr]=\sum_{m=-1}^3\pi_m\,\Res_Y\bigl[Y^m\re^{-W_V/\zh}\,\dd Y\bigr],
\]
the products being formed in the hatted variables, with finite fibres, because $\pi_m$ is a polynomial in $\zh$ at each $u$-order. The model residues are, term by term,
\[
\Res_{Y=0}\bigl[Y^m\re^{-W_V/\zh}\,\dd Y\bigr]=\sum_{i'\geq0}\frac{(-1/(4\zh))^{i'}}{i'!}\,\frac{(-V/\zh)^e}{e!}=s^{-(m+1)}\IC_m(v(u),z),\qquad e=m+4i'+1:
\]
the substitution $p=sY$, with $v=uV$ and $z=s^4\zh$, turns $\Wv(p)/z$ into $W_V(Y)/\zh$, and expanding $V^e$ in $u$ is a relabelling of the kind treated in \cref{lem:con-relab}. The exponent map $\iota$ of \cref{lem:con-relab} is additive, so the image of a product with finite fibres is the product of the images. Hence $\cV=s\sum_m\pi_m\,s^{-(m+1)}\IC_m=\Sigma$, the formal sum of \cref{lem:con-key}, coefficient by coefficient. By \cref{lem:con-id}, $\pi_{\nu_b}=\varepsilon_b\varpi_0G_b=\varepsilon_b\varpi_0s^{\nu_b}P_{0b}$, and $\YC_b=\varepsilon_b\IC_{\nu_b}$ (\cref{def:gau-van}); the signs cancel, $\varepsilon_b^2=1$, and so do the powers of $s$, $s\cdot s^{-(\nu_b+1)}\cdot s^{\nu_b}=1$; so $\Sigma=\varpi_0\sum_bP_{0b}\YC_b=\varpi_0(P\YC)_0$. No power of $\tau$ survives: the residue side carries none, and $G_b$ is defined without it. Since $\cV\in z^{-1}\bbQ[z^{-1}]\fps{u}$ (\cref{def:con-PhiSP}), the identity holds in $\bbQ\lser{z}\fps{u}$.
\end{proof}

\begin{thm}[Theorem P(b)]\genusstatus{AI-proposed, unchecked}
\label{thm:con-Pb}
In every conifold frame of \cref{def:set-frames}, $\PhiSP=\Phican$ in $\bbQ\fps{u}\fps{\zh}$.
\end{thm}

Besides the lemmas cited in its steps, the proof uses Theorem~U (\cref{thm:gau-U}) for the pair of systems of \cref{prop:gau-quintic,prop:gau-equalprod}, whose hypotheses hold in the conifold frames, the gauge identity \cref{thm:gau-G}, the normalisations $\varpi_0(0)=1$ and $\tc'(0)=1$ of the conifold frames, and $\mathrm{ph}^{\mathrm{can}}_0=1$ (\cref{lem:gau-diag}, through \cref{thm:gau-R0}~(ii)). It does not use (QRR).

\begin{proof}
Steps 1--4 are \cref{lem:con-cfu} (the uniformising map), \cref{lem:con-red} (the reduction modulo exact forms), the invariance of both functionals under the uniformising map and their vanishing on exact forms (\cref{lem:fsp-cfu,lem:fsp-ibp,lem:fsp-res}), and the identification of \cref{lem:con-S2,lem:con-id}.

\emph{Step 5 (regularity).} By \cref{thm:con-Pa}~(iii), $\PhiSP$ solves HPF; by \cref{thm:gau-R0}~(ii), which rests on \cref{cor:gau-cycsys}, so does $\Phican$; both have constant term $1$ and coefficients in $\bbQ\fps{u}$. By \cref{lem:gau-poly}, $\PhiSP=\Phican\,\mathsf s(z)$ for a series $\mathsf s\in1+z^5\bbQ\fps{z^5}$ (the letter $\mathsf s$ is local to this proof; $z^5=u^4\zh^5$). Put $P^{\mathrm{SP}}:=\mathsf s(z)\Pc$. By \cref{thm:gau-U}~(b) and \cref{prop:gau-graded} it is a graded regular gauge, and $\Phi_{P^{\mathrm{SP}}}=\mathsf s(z)\Phican=\PhiSP$. The gauge identity for the canonical gauge, $\Pc\YC(v(u),z)=\YCQ$ (\cref{thm:gau-G}), multiplied by $\mathsf s(z)$, gives $(P^{\mathrm{SP}}\YC)_0=\mathsf s(z)(\YCQ)_0$.

\emph{Step 6 (the residue).} By \cref{prop:con-res} for $P=P^{\mathrm{SP}}$, $\cV=\varpi_0\,\mathsf s(z)\,(\YCQ)_0$ in $\bbQ\lser{z}\fps{u}$.

\emph{Step 7 (the order $u^1$).} On the left only $k=1$ contributes to $u^1$ in \cref{def:con-PhiSP}, because $\Xi=u+O(u^2)$, and $[x^0](\re^{-\cG/z}\fa)=1$, because $\cG=O(x^4)$ and $\fa(0,z)=1$; so $[u^1]\cV=-1/z$. On the right, $[u^0](\YCQ)_0=0$, since $y^{(0)}=-e_1$ and the terms $y^{(n)}$, $n\geq1$, vanish at $u=0$; and $[u^1](\YCQ)_0=-1/z$, since by \cref{lem:gau-vanq} only $y^{(1)}_0=-\tc$ contributes at $u^1$ and $\tc=u+O(u^2)$. With $\varpi_0(0)=1$ the right side has $u^1$-coefficient $-\mathsf s(z)/z$. Hence $\mathsf s=1$, and $\PhiSP=\Phican$.
\end{proof}

% Section 8 (part 8): the glue. Last part of Section 8.

\subsection{The glue}
\label{sec:con-glue}

For a colour-$0$ WKB function with polynomial coefficients write $\Phi(z,1):=\sum_n\mathrm{ph}_n(1)z^n$, its value at $u=1$, where $\zh=z$.

\begin{thm}[The glue]\genusstatus{AI-proposed, unchecked}
\label{thm:con-glue}
In $\bbQ(\zeta)\fps{z}$,
\begin{equation}
\label{eq:con-glue}
h(z)\,\Dn_0(z)=\Phican(z,1);
\end{equation}
hence $\Dn_0\in1+z\bbQ\fps{z}$, and, by \cref{thm:con-Pb}, $\Dn_0(z)=\PhiSP(z,1)/h(z)$.
\end{thm}

\begin{proof}
(a) By \cref{thm:gau-R0}~(ii)--(iv), the canonical flat amplitude $A^{\mathrm{can}}=\sum_nA^{\mathrm{can}}_nz^n$ has $A^{\mathrm{can}}_n=L^{4n+1}\mathrm{ph}^{\mathrm{can}}_n(L^{-5})\in\bbQ[L,L^{-1}]$, the $\mathrm{ph}^{\mathrm{can}}_n$ being polynomials. In $\bbQ(\zeta)\fps{q}$, with $L=(1-5^5q)^{-1/5}$, $A^{\mathrm{can}}$ and the actual amplitude $A^{\mathrm{act}}$ are colour-$0$ WKB columns of $\PFfive$ with the same exponential $\re^{u_0/z}$ and the same leading term $L$, so by formal WKB uniqueness (\cref{lem:gau-wkb}, with $\Bbbk=\bbQ(\zeta)$) they differ by a constant series in $1+z\bbQ(\zeta)\fps{z}$; hence every $A^{\mathrm{act}}_n$ lies in $\bbQ(\zeta)[L,L^{-1}]$. The ring $\bbQ(\zeta)[L,L^{-1}]\fps{z}$ embeds in $\bbQ(\zeta)\lser{s}\fps{\zh}$, where the comparison of \cref{thm:gau-R0} reads $A^{\mathrm{act}}\,\Dn_0=A^{\mathrm{can}}$; so the constant series is $\Dn_0$. (b) Evaluation at $L=1$, that is $u=1$, $q=0$, $\zh=z$, is a ring homomorphism $\bbQ(\zeta)[L,L^{-1}]\to\bbQ(\zeta)$; applied coefficientwise it gives $A^{\mathrm{act}}(z)|_{L=1}\,\Dn_0(z)=\Phican(z,1)$. As $(1-5^5q)^{-1/5}=1+O(q)$, $A^{\mathrm{act}}_n|_{L=1}$ is the coefficient of $q^0$ of $A^{\mathrm{act}}_n$, which is $h_n$ by (QRR) in the form of \cref{thm:set-qrr}, \eqref{eq:set-qrr}. As $h\in1+z\bbQ\fps{z}$ is invertible, $\Dn_0=\Phican(z,1)/h(z)$ has rational coefficients, and \cref{thm:con-Pb} gives the second form. The inputs used here are the value at $q=0$ of (QRR) and the actual rows with their WKB form, which rest on (M), (G0) (\cref{thm:set-mirror,thm:set-G0}) and the flatness part of (QRR) (\cref{def:set-rows}~(b)) through \cref{cor:set-cyclicPF5} and \cref{thm:gau-R0}; no other step of this section uses an input.
\end{proof}

\begin{cor}\genusstatus{AI-proposed, unchecked}
\label{cor:con-A2}
$\Dn=1$ if and only if $\PhiSP(z,1)=h(z)$. By \cref{thm:con-Pa}~(iv), $\PhiSP(z,1)=c(z;1)$, the value at $L=1$ of the Laurent polynomials $c_n$; so the condition reads $c(z;1)=h(z)$ (\cref{cor:gam-equiv}).
\end{cor}

\begin{proof}
By \cref{cor:gau-colours}, $\Dn_i(z)=\Dn_0(\zeta^{-i}z)$, so $\Dn=1$ if and only if $\Dn_0=1$; by \eqref{eq:con-glue} and \cref{thm:con-Pb}, $\Dn_0=1$ if and only if $\PhiSP(z,1)=h(z)$. By \cref{thm:con-Pa}~(iv), $c_n(1)=\mathrm{ph}^{\mathrm{SP}}_n(1)$.
\end{proof}

